% Source-aligned English for printed p.189.
\clearpage
\SourcePage{189}

\LayerTarget{tohoku:eq:4.2.8}
\LayerTarget{tohoku:proof:theorem-4.2.2}
With $A$ fixed, these homomorphisms are isomorphisms for every $p$ and $B$ if
and only if (i) this is so for $p=0$, and (ii)
$\mathrm{Ext}_{O(x)}^p(A(x),B(x))=0$ for $p>0$ whenever $B$ is an injective
$O$-module. Propositions 4.1.1 and 4.1.2 were plainly designed to verify these
conditions. We therefore obtain:

\LayerTarget{tohoku:theorem:4.2.2}
\noindent\textsc{Theorem 4.2.2.} \textit{Suppose that $O$ is a coherent sheaf
of left Noetherian rings and that $A$ is a coherent $O$-module. Then the
homomorphisms (4.2.8) are isomorphisms for every $B$ and $p$.}

\LayerTarget{tohoku:remark:locally-free-ext-vanishing}
Finally, note the following trivial case in which the homomorphisms (4.2.8)
are isomorphisms because both sides vanish:

\LayerTarget{tohoku:proposition:4.2.3}
\noindent\textsc{Proposition 4.2.3.} \textit{Suppose that $A$ is locally
isomorphic to $O^n$, where $n>0$ is a given finite integer. Then
$\mathbf{Ext}_O^p(A,B)=0$ for every $O$-module $B$ and every $p>0$.}

\LayerTarget{tohoku:proof:proposition-4.2.3}
Indeed, by the local nature of $\mathbf{Ext}_O^p(A,B)$, one may suppose that
$A=O^n$. Then $\mathbf{Hom}_O(A,B)=B^n$ is an exact functor in $B$, which
proves the assertion.

\LayerTarget{tohoku:corollary:4.2.3}
\noindent\textsc{Corollary.} \textit{If $A$ is locally isomorphic to $O^n$,
then}
\[
\mathrm{Ext}_O^p(X;A,B)=H^p(X,\mathbf{Hom}_O(A,B))
\]
\textit{for every $O$-module $B$.}

\LayerTarget{tohoku:proof:corollary-4.2.3}
This follows at once from the spectral sequence of Theorem 4.2.1. In
particular, taking $A=O$, one obtains

\LayerTarget{tohoku:eq:4.2.9}
\[
\tag{4.2.9}
H^p(X,B)=\mathrm{Ext}_O^p(X;O,B).
\]

\LayerTarget{tohoku:remark:cohomology-as-ext}
This formula is valid without any hypothesis on $O$ or on the $O$-module $B$;
it is also immediate directly, since $H^p(X,B)$ are the derived functors,
\textit{in $C^O$} (by Proposition 4.1.3), of
$\Gamma(B)=\mathrm{Hom}_O(O,B)$.

\LayerTarget{tohoku:construction:computing-spectral-sequence}
To conclude this subsection, we indicate a method for computing the spectral
sequence of Theorem 4.2.1 which is more convenient than a direct application
of the definition.

\LayerTarget{tohoku:proposition:4.2.4}
\noindent\textsc{Proposition 4.2.4.} \textit{Let
$\mathbf{L}=(L_i)$ be a left resolution of $A$ by $O$-modules $L_i$ locally
isomorphic to modules $O^{n_i}$. Consider the functor
$\mathbf{h}_{\mathbf{L}}$ on $C^O$, with values in complexes in nonnegative
degrees in $C$, defined by
$\mathbf{h}_{\mathbf{L}}(B)=\mathbf{Hom}_O(\mathbf{L},B)$, and likewise put
$\mathbf{h}_A(B)=\mathbf{Hom}_O(A,B)$. Then $\mathbf{h}_{\mathbf{L}}$ is a
resolvent functor (cf. 2.5) for $\mathbf{h}_A$.}

\LayerTarget{tohoku:proof:proposition-4.2.4}
Indeed, $\mathbf{h}_{\mathbf{L}}$ is exact by Proposition 4.2.3. Moreover,
$H^0(\mathbf{h}_{\mathbf{L}}(B))=\mathbf{h}_A(B)$ because the functor
$C\longmapsto\mathbf{Hom}_O(C,B)$ is left exact in $C$. Finally, if $B$ is
injective, then $\mathbf{h}_{\mathbf{L}}(B)$ is acyclic in dimensions $>0$,
since in that case $C\longmapsto\mathbf{Hom}_O(C,B)$ is exact.

\LayerTarget{tohoku:corollary:4.2.4:1}
\noindent\textsc{Corollary 1.} \textit{Under the hypotheses of the preceding
proposition, one has}

\TohokuContinuation{190}
